\documentclass[a4paper,11pt]{article}
\usepackage[T1]{fontenc}
\usepackage[utf8]{inputenc}
\usepackage{t1enc}
\usepackage[dvips]{color}
\usepackage{amsmath,amssymb,makeidx,stmaryrd}
\usepackage[normalem]{ulem}
\usepackage{fancybox}
\definecolor{gris}{gray}{0.9}
\usepackage{pst-plot,pst-tree,pstricks,pst-node,pst-text}
\usepackage{pst-eucl}
\usepackage{pstricks-add}
\usepackage[left=3.5cm, right=3.5cm, top=3cm, bottom=3cm]{geometry}
\newcommand{\vect}[1]{\overrightarrow{\,\mathstrut#1\,}}
\def\bas{\fcolorbox{gris}{gris}{\parbox{0.986\textwidth}{
	\color{red}\footnotesize 1994\hfill Page : \thepage/1}}}
\def\haut{\fcolorbox{gris}{gris}{\parbox{0.986\textwidth}{
	\color{red}
	\bf Concours Général de Mathématiques \hfill 1994
	}}}
\makeatletter
\def\ps@cg{%
\let\@oddhead\haut%
\let\@evenhead\haut%
\let\@oddfoot\bas
\let\@evenfoot\bas}
\makeatother
\newcommand{\R}{\mathbb{R}}
\newcommand{\N}{\mathbb{N}}
\newcommand{\D}{\mathbb{D}}
\newcommand{\Z}{\mathbb{Z}}
\newcommand{\Q}{\mathbb{Q}}
\newcommand{\C}{\mathbb{C}}
\pagestyle{cg}
\renewcommand{\theenumi}{\textbf{\arabic{enumi}}}
\renewcommand{\labelenumi}{\textbf{\theenumi.}}
\renewcommand{\theenumii}{\textbf{\alph{enumii}}}
\renewcommand{\labelenumii}{\textbf{\theenumii.}}
\parindent0pt
\everymath{\displaystyle}
\usepackage[french]{babel}

\DecimalMathComma
\usepackage[np]{numprint}
\begin{document}

\begin{center}\textbf{{\color{blue} Exercice I}}
\end{center}

Pour tout entier naturel $n$, on note $I_n$ le nombre d'entiers $p$ pour lesquels

\[50^n < 7^p < 50^{n+1}\]

\smallskip

\begin{enumerate}
\item Démontrer que pour tout entier $n$, \:$I_n$ vaut 2 ou 3.
\item Démontrer qu'il existe une infinité d'entiers $n$ pour lesquels $I_n$ vaut 3 et donner le plus petit d'entre eux.
\end{enumerate}

\begin{center}\textbf{{\color{blue} Exercice II}}
\end{center}

Soit $\sum$ une demi-sphère et $P$ le plan contenant son cercle de base. Un plan variable $Q$, parallèle à un plan fixe non perpendiculaire à $P$, coupe $\sum$ suivant un cercle $C$. On désigne par $C'$ le projeté orthogonal de $C$ sur $P$.

Comment doit-on placer le plan $Q$ pour que le cylindre de bases $C$ et $C'$ ait un volume maximal ?

\begin{center}\textbf{{\color{blue} Exercice III}}
\end{center}

On définit une application $f$ de $\N$ dans $\N$ par $f(1) = 0$ et pour tout élément $n$ de $\N^{*}$  :

\[f(2n) = 2f(n) + 1 \qquad  f(2n+1) = 2f(n)\]

Etant donné un entier strictement positif $p$ quelconque, on pose $u_0 = p$ et, tant que $u_k$ appartient à $\N^{*}$

\[u_{k+1} = f\left(u_k\right)\]

\medskip

\begin{enumerate}
\item Montrer que, pour tout choix de $p$, il existe un unique entier $v(p)$ tel que $u_{v(p)} = 0$.
\item 
	\begin{enumerate}
		\item Calculer $v(\np{1994})$ .
		
Quel est le plus petit entier $p$ tel que $v(p) = v(\np{1994})$ ?
		\item Étant donné un entier $N$, déterminer le plus petit entier $p$ tel que $v(p) = N$.
	\end{enumerate}
\end{enumerate}

\begin{center}\textbf{{\color{blue} Exercice IV}}
\end{center}

Soit ABC un triangle. Si P est un point de son plan, on note L, M, N les projetés orthogonaux de P respectivement sur les droites (BC), (CA) et (AB). 

Déterminer le point P pour lequel la quantité 

\[\text{BL}^2  + \text{CM}^2 + \text{AN}^2\]

est minimale.


\begin{center}\textbf{{\color{blue} Exercice V}}
\end{center}

Soit $f$ une application de $\N$ dans $\N$ telle que $f(1) > 0$ et, quels que soient les entiers naturels $m$ et $n$, on a :

\[f\left(m^2 + n^2\right) = [f(m)]^2 + [f(n)]^2\]

\begin{enumerate}
\item Calculer $f (k)$ pour $0 \leqslant k \leqslant 12$.
\item Calculer $f(n)$ , $n$ étant un entier quelconque.
\end{enumerate}
\end{document}